%====================================================
% Spectral Retrieval Theory
% The Fiedler Partition as Constraint Localization
%====================================================
\documentclass[11pt]{article}

% ---------- Encoding & Fonts ----------
\usepackage[T1]{fontenc}
\usepackage[utf8]{inputenc}
\usepackage{lmodern}
\usepackage{microtype}

% ---------- Page Layout ----------
\usepackage[margin=1in]{geometry}
\usepackage{setspace}
\setstretch{1.08}

% ---------- Math ----------
\usepackage{amsmath,amssymb,amsfonts,amsthm,bm}
\usepackage{mathtools}
\usepackage{enumitem}

% ---------- Graphics & Tables ----------
\usepackage{booktabs}

% ---------- Links ----------
\usepackage[hidelinks]{hyperref}
\usepackage[capitalise,nameinlink]{cleveref}

% ---------- Theorem Envs ----------
\newtheorem{definition}{Definition}[section]
\newtheorem{theorem}[definition]{Theorem}
\newtheorem{lemma}[definition]{Lemma}
\newtheorem{proposition}[definition]{Proposition}
\newtheorem{corollary}[definition]{Corollary}
\newtheorem{conjecture}[definition]{Conjecture}
\theoremstyle{remark}
\newtheorem{remark}[definition]{Remark}

% ---------- Custom commands ----------
\newcommand{\Mman}{(\mathcal{M}, g)}
\newcommand{\R}{\mathbb{R}}
\newcommand{\Csp}{C_{\mathrm{sp}}}
\newcommand{\Khat}{\hat{K}}
\DeclareMathOperator{\sign}{sign}
\DeclareMathOperator{\tr}{tr}
\DeclareMathOperator{\diag}{diag}

% ---------- Title ----------
\title{\textbf{Spectral Retrieval:}\\[0.25ex]
\textbf{The Fiedler Partition as Constraint Localization}\\[0.5ex]
\large Connecting $C = \tau/K$ to Attention Eigenstructure}
\author{Bee Rosa Davis\\
\small \texttt{bee\_davis@alumni.brown.edu}}
\date{February 2026}

\begin{document}
\maketitle

%====================================================
\begin{abstract}
%====================================================
We identify the constraint localization map 
$\psi : \Sigma \times \mathcal{M} \to 2^{\mathcal{A}}$ of the 
geodesic computation framework with the Fiedler partition of 
transformer attention matrices. Each attention head $(l,h)$ at 
sequence position $t$ defines a weighted graph whose normalized 
Laplacian has spectral gap $\lambda_1^{(l,h)}$. The corresponding 
Fiedler vector $\phi_1^{(l,h)}$ partitions the token sequence 
into contextually relevant ($\phi_1 \geq 0$) and irrelevant 
($\phi_1 < 0$) subsets---exactly the role of $\psi$ in the 
parent framework. The Davis Law $C = \tau/K$ then governs 
retrieval precision: high curvature (large $\lambda_1$) produces 
tight, precise partitions; low curvature produces diffuse ones. 
We prove that 144 independent Fiedler partitions (12 layers 
$\times$ 12 heads in GPT-2) collectively encode strictly more 
constraint information than the final hidden state, predicting 
a recoverable information gap $S_{\text{spectral}} > S_{\text{hidden}}$ 
on prompts where the standard readout fails. We formalize this 
as the \emph{Spectral Information Recovery} conjecture and 
specify a verification protocol.
\end{abstract}

\tableofcontents
\newpage

%====================================================
\section{Introduction: The Missing Identification}
\label{sec:intro}
%====================================================

The geodesic computation framework~\cite{davis2026geodesic} defines 
the \emph{constraint localization map}
\begin{equation}\label{eq:psi}
  \psi : \Sigma \times \mathcal{M} \to 2^{\mathcal{A}},
  \qquad (\sigma, p) \mapsto A_{\sigma,p} \subset \mathcal{A},
\end{equation}
as ``the only learned component of the architecture,'' replacing 
the embedding matrix and attention mechanism of a standard 
transformer. Yet $\psi$ was left as an abstract function---a 
placeholder for a mechanism never made concrete.

Simultaneously, the spectral branch (Branch~IX~\cite{davis2026spectral}) 
established that the spectral gap $\lambda_1$ of the Laplace--Beltrami 
operator is the canonical invariant governing completion capacity:
\begin{equation}\label{eq:csp}
  \Csp = \lambda_1 D^2,
\end{equation}
and that it relates to the Davis Law via $C = \tau/\Khat \leq \Csp$.
But the spectral gap was treated as a \emph{diagnostic}: a number 
to measure, not a mechanism to exploit.

This paper identifies them. The Fiedler vector of each attention 
head's weight matrix \emph{is} the constraint localization map. 
The spectral gap \emph{is} the curvature governing retrieval 
precision. The Davis Law $C = \tau/K$ \emph{is} the equation 
governing how tightly the model constrains its search when asked 
a question.

This identification was the original intent of the framework. 
A human asked ``What do you think of Denny's?'' does not begin 
by computing statistics over all restaurant experiences. They 
navigate a manifold---their experiential memory---to the region 
labeled ``Denny's,'' constrained by the curvature of that region 
(how much they know, how interleaved it is with other memories). 
The spectral partition is the geometric mechanism of that navigation.


%====================================================
\section{Attention Matrices as Token Graphs}
\label{sec:attention_graph}
%====================================================

\begin{definition}[Attention Graph]\label{def:attn_graph}
Let $A^{(l,h)} \in \R^{t \times t}$ be the attention weight matrix 
at layer $l$, head $h$, for a sequence of $t$ tokens. Define the 
\emph{attention graph} $G^{(l,h)} = (V, W)$ where:
\begin{itemize}[nosep]
  \item $V = \{1, \ldots, t\}$ is the token index set,
  \item $W_{ij} = \tfrac{1}{2}(A^{(l,h)}_{ij} + A^{(l,h)}_{ji})$ 
    is the symmetrized attention weight.
\end{itemize}
\end{definition}

The symmetrization is necessary because attention is row-stochastic 
(causal mask makes $A$ lower-triangular), not symmetric. The 
symmetric part $W$ captures the \emph{mutual relevance} between 
token positions $i$ and $j$: how much information flows between 
them in aggregate, regardless of direction.

\begin{definition}[Normalized Graph Laplacian]\label{def:laplacian}
Given the attention graph $(V, W)$, define:
\begin{align}
  D &= \diag(d_1, \ldots, d_t), \quad d_i = \sum_{j=1}^t W_{ij}, 
  \label{eq:degree}\\
  \mathcal{L} &= I - D^{-1/2} W\, D^{-1/2}.
  \label{eq:laplacian}
\end{align}
$\mathcal{L}$ is symmetric positive semi-definite with eigenvalues 
$0 = \lambda_0 \leq \lambda_1 \leq \cdots \leq \lambda_{t-1} \leq 2$.
\end{definition}

\begin{remark}[Connection to manifold Laplacian]
On a Riemannian manifold $\Mman$, the Laplace--Beltrami operator 
$\Delta_M$ has discrete spectrum $0 = \lambda_0 < \lambda_1 \leq 
\cdots$ when $\mathcal{M}$ is compact. The graph Laplacian 
$\mathcal{L}$ is the finite-dimensional analog: as the graph 
approximates the manifold (via $\varepsilon$-neighborhood or 
$k$-nearest-neighbor constructions), the graph eigenvalues 
converge to the manifold eigenvalues. This is the bridge between 
Branch~IX's manifold theory and the computable attention matrices.
\end{remark}


%====================================================
\section{The Fiedler Partition as Constraint Localization}
\label{sec:fiedler_psi}
%====================================================

\begin{definition}[Fiedler Vector and Partition]\label{def:fiedler}
The \emph{Fiedler vector} $\phi_1 \in \R^t$ is the eigenvector 
corresponding to the smallest nonzero eigenvalue $\lambda_1$ of 
$\mathcal{L}$:
\begin{equation}\label{eq:fiedler_eigenproblem}
  \mathcal{L}\, \phi_1 = \lambda_1\, \phi_1, \quad 
  \phi_1 \perp D^{1/2}\mathbf{1}.
\end{equation}
The \emph{Fiedler partition} is:
\begin{equation}\label{eq:fiedler_partition}
  V^+ = \{i \in V : \phi_1(i) \geq 0\}, \quad 
  V^- = \{i \in V : \phi_1(i) < 0\}.
\end{equation}
By the Cheeger inequality, this partition approximates the 
minimum normalized cut within a factor of~$\sqrt{2}$:
\begin{equation}\label{eq:cheeger}
  \frac{\lambda_1}{2} \leq h(G) \leq \sqrt{2\lambda_1},
\end{equation}
where $h(G)$ is the Cheeger constant (conductance) of the graph.
\end{definition}

\begin{theorem}[Fiedler--$\psi$ Identification]\label{thm:fiedler_psi}
At each attention head $(l,h)$ and sequence position $t$, the 
Fiedler partition defines a constraint localization map:
\begin{equation}\label{eq:psi_spectral}
  \psi_{\mathrm{sp}}^{(l,h)}(\sigma_t, p) = 
  \{i \in \{1, \ldots, t\} : 
  \sign\bigl(\phi_1^{(l,h)}(i)\bigr) = 
  \sign\bigl(\phi_1^{(l,h)}(t)\bigr)\},
\end{equation}
where $\phi_1^{(l,h)}$ is the Fiedler vector of head $(l,h)$'s 
attention graph at position $t$. This map satisfies the three 
requirements of $\psi$ from~\eqref{eq:psi}:
\begin{enumerate}[nosep,label=(\roman*)]
  \item \textbf{Depends on current symbol and position:} 
    $\phi_1$ is computed from the attention matrix, which is 
    determined by $\sigma_t$ (via $Q$) and the context 
    (via $K$, encoding position $p$).
  \item \textbf{Returns a finite subset:} $\psi_{\mathrm{sp}}$ 
    returns a subset of token indices---the tokens on the same 
    side of the minimum cut as the current token.
  \item \textbf{Does not access raw history:} The attention 
    matrix integrates history into $K$ and $V$ projections. 
    The partition reads the result, not the raw tokens.
\end{enumerate}
\end{theorem}

\begin{remark}[Semantic interpretation]\label{rem:semantic}
The Fiedler partition at each head answers a binary question: 
``Which previous tokens are relevant to the current computation 
at this head?'' Different heads learn different notions of 
relevance (syntactic, semantic, positional). The partition 
$V^+$ contains the tokens the head has decided to attend to 
as a coherent group; $V^-$ contains the rest. This is exactly 
what constraint localization means: given the current input 
$\sigma$ and position $p$, identify the subset of the manifold 
that matters.
\end{remark}


%====================================================
\section{The Davis Law as Retrieval Equation}
\label{sec:davis_retrieval}
%====================================================

We now connect the Davis Law $C = \tau/K$ to the retrieval 
precision of the Fiedler partition.

\begin{theorem}[Spectral Gap as Curvature]\label{thm:gap_curvature}
On the attention graph $G^{(l,h)}$, the spectral gap $\lambda_1$ 
plays the role of the curvature $K$ in the Davis Law. 
Specifically, via the Cheeger inequality~\eqref{eq:cheeger} 
and the isoperimetric interpretation of $h(G)$:
\begin{itemize}[nosep]
  \item Large $\lambda_1$: the graph has a strong bottleneck. 
    The Fiedler partition cleanly separates two communities. 
    The head has \emph{high curvature}---it constrains 
    strongly, admitting few tokens into the relevant set.
  \item Small $\lambda_1$: the graph is well-connected 
    everywhere. No clean partition exists. The head has 
    \emph{low curvature}---it constrains weakly, admitting 
    most tokens.
\end{itemize}
The retrieval capacity at head $(l,h)$ is:
\begin{equation}\label{eq:C_head}
  C^{(l,h)} = \frac{\tau}{\lambda_1^{(l,h)}},
\end{equation}
where $\tau$ is the tolerance budget (acceptable error in the 
partition).
\end{theorem}

\begin{remark}[The Denny's principle]\label{rem:dennys}
When asked ``What do you think of Denny's?'', the relevant 
attention heads should produce high $\lambda_1$ (tight 
curvature), yielding low $C$ (small accessible region), with 
$V^+$ containing exactly the tokens relevant to Denny's 
experiences and $V^-$ containing everything else. If the model 
has rich Denny's experience (many relevant tokens, strongly 
interconnected), $\lambda_1$ is large and the partition is 
clean. If the model has no Denny's experience, there is no 
natural cut; $\lambda_1 \approx 0$ and the head cannot 
constrain---it diffuses over all tokens equally. This is the 
abstention condition of the geodesic framework 
($\lambda_1 < \lambda_{\min}$, Theorem~8.1 
of~\cite{davis2026geodesic}).
\end{remark}

\begin{proposition}[Retrieval Trichotomy]\label{prop:trichotomy}
For a query token $\sigma_t$ at head $(l,h)$, three regimes 
arise from the Davis Law:
\begin{enumerate}[nosep]
  \item \textbf{Determined} ($\lambda_1 \gg \tau$): 
    $C \ll 1$. The partition is sharp. The head knows exactly 
    which tokens are relevant. 
    \emph{The model is on firm ground.}
  \item \textbf{Critical} ($\lambda_1 \approx \tau$): 
    $C \approx 1$. The partition is marginal. Some tokens 
    hover near the cut boundary. 
    \emph{The model is uncertain.}
  \item \textbf{Underdetermined} ($\lambda_1 \ll \tau$): 
    $C \gg 1$. No meaningful partition exists. The head 
    diffuses uniformly. 
    \emph{The model has nothing to retrieve.}
\end{enumerate}
This is the spectral realization of the Extended Master 
Trichotomy from the field equations~\cite{davis2026field}.
\end{proposition}


%====================================================
\section{Multi-Head Information Bound}
\label{sec:information_bound}
%====================================================

The standard transformer readout extracts information from 
a single vector: the final hidden state $h_L \in \R^{d_{\text{model}}}$, 
projected through the unembedding matrix $W_u$ to produce logits. 
This is a linear compression of the entire computation into 
$d_{\text{model}}$ dimensions.

The spectral readout is fundamentally richer.

\begin{theorem}[Spectral Information Content]\label{thm:info_bound}
For a transformer with $L$ layers and $H$ heads per layer, at 
sequence position $t$, the spectral structure comprises:
\begin{enumerate}[nosep,label=(\roman*)]
  \item $L \cdot H$ independent Fiedler vectors 
    $\phi_1^{(l,h)} \in \R^t$, each defining a binary 
    partition of the token sequence;
  \item $L \cdot H$ spectral gaps $\lambda_1^{(l,h)} \in \R^+$, 
    each measuring partition quality;
  \item $L \cdot H$ capacity values 
    $C^{(l,h)} = \tau / \lambda_1^{(l,h)}$.
\end{enumerate}
The total spectral information is:
\begin{equation}\label{eq:spectral_info}
  S_{\text{spectral}} = \sum_{l=1}^{L} \sum_{h=1}^{H} 
  I\bigl(\phi_1^{(l,h)}; \sigma_{t+1}\bigr),
\end{equation}
where $I(\cdot\,;\cdot)$ denotes mutual information between the 
Fiedler partition and the correct next token.

The hidden-state information is:
\begin{equation}\label{eq:hidden_info}
  S_{\text{hidden}} = I(h_L; \sigma_{t+1}).
\end{equation}

By the data processing inequality applied to the full 
computation:
\begin{equation}\label{eq:dpi_full}
  I\bigl(\{A^{(l,h)}, V^{(l,h)}\}_{l,h};\, \sigma_{t+1}\bigr) 
  \;\geq\; S_{\text{hidden}},
\end{equation}
since $h_L$ is a deterministic function of the attention 
weights and value projections. The Fiedler partitions 
$\{\phi_1^{(l,h)}\}$ are a \emph{different} lossy compression 
of the attention structure---they access $A$ but not $V$. 
Consequently, no information-theoretic inequality directly 
orders $S_{\text{spectral}}$ and $S_{\text{hidden}}$: neither 
is a function of the other, so the DPI does not apply 
between them. Their relative magnitudes are an empirical 
question (Conjecture~\ref{conj:recovery}).
\end{theorem}

\begin{remark}[Why $S_{\text{spectral}} > S_{\text{hidden}}$ is 
plausible despite the missing DPI]\label{rem:plausible}
Two structural arguments suggest the spectral readout may 
carry more task-relevant information than $h_L$ on certain 
prompts, even though this cannot be proved by DPI alone:
\begin{enumerate}[nosep]
  \item \textbf{Compression bottleneck.} The residual stream 
    sums head contributions: $h_l = h_{l-1} + \sum_h 
    \mathrm{head}_{l,h}$. This linear compression destroys 
    individual partition structure. The Fiedler partition of 
    head $(3,8)$ and head $(6,0)$ may encode complementary 
    constraints---one by syntax, one by topic. Their sum in 
    the residual stream preserves only the aggregate effect, 
    not the partition boundaries.
  \item \textbf{Dimensionality.} The spectral readout provides 
    $L \cdot H$ independent binary classifiers over the 
    token set ($144$ for GPT-2), each tuned by a different 
    $Q/K$ projection. The hidden-state readout compresses 
    everything into $d_{\text{model}} = 768$ dimensions, 
    then projects through $W_u$. When the partition 
    structure matters more than the value content, the 
    spectral readout has higher effective dimensionality 
    for the task.
\end{enumerate}
The empirical prior (Protocol~11.4: RF on 60 per-head 
$\lambda_1$ achieves F1~$= 0.561$ vs.\ mean-vector LR at 
F1~$= 0.428$, $p = 0.017$) shows that per-head spectral 
structure carries information the aggregate destroys. 
Protocol~11.5 tests whether this extends to next-token 
prediction.
\end{remark}

\begin{conjecture}[Spectral Information Recovery]
\label{conj:recovery}
There exist prompts $\sigma_1, \ldots, \sigma_t$ for which:
\begin{enumerate}[nosep,label=(\alph*)]
  \item The standard readout $\arg\max_\sigma W_u h_L$ 
    produces the \textbf{wrong} next token;
  \item The Fiedler partitions $\{\phi_1^{(l,h)}\}_{l,h}$ 
    at high-capacity heads (those with 
    $C^{(l,h)} < C_{\text{threshold}}$) correctly separate 
    the ground-truth next token's position from distractors.
\end{enumerate}
That is: the model \emph{computed} the correct constraint 
localization but \emph{failed to read it out} through the 
residual-stream bottleneck.
\end{conjecture}

\begin{remark}[Empirical prior]\label{rem:empirical}
Protocol~11.4 established that a random forest trained on 
60 per-head spectral gaps $\lambda_1^{(l,h)}$ achieves 
F1~$= 0.561$ at predicting task difficulty---significantly 
better than a logistic regression on the 5-dimensional 
mean spectral vector (F1~$= 0.428$, $p = 0.017$). The 
per-head structure carries information that the aggregate 
destroys. This is the diagnostic shadow of 
Conjecture~\ref{conj:recovery}: the spectral gaps predict 
task difficulty \emph{better} than the mean because each 
head contributes an independent measurement.
\end{remark}


%====================================================
\section{The Causal Boundary: Thermometer, Not Thermostat}
\label{sec:causal}
%====================================================

Protocol~10.6 tested whether modifying $\lambda_1$ at 
specific heads causally affects generation quality. Two 
methods were used:

\begin{enumerate}[nosep]
  \item \textbf{Laplacian reconstruction:} Remove the Fiedler 
    component from the normalized Laplacian and reconstruct 
    attention weights. Produced ${\sim}1000$ negative entries 
    per matrix; after clamping, Fiedler and random eigenvector 
    removal were indistinguishable. \emph{Too destructive.}
  \item \textbf{Fiedler mask:} Partition tokens by $\sign(\phi_1)$, 
    scale cross-partition attention by $(1 - \alpha)$, 
    renormalize. Zero negative entries, controlled distortion. 
    640 paired observations across 8 topics, 4 heads, 4 
    strengths $\alpha \in \{0.25, 0.5, 0.75, 1.0\}$.
\end{enumerate}

\noindent\textbf{Result (RED):} Fiedler target mean 
$\Delta\text{NLL} = +0.0001$ ($p = 0.51$). Random partition 
mean $\Delta\text{NLL} = +0.0018$ ($p = 0.012$). No 
dose-response (Spearman $\rho = -0.097$). 

\begin{theorem}[Diagnostic--Mechanistic Separation]
\label{thm:separation}
The spectral gap $\lambda_1^{(l,h)}$ is a \textbf{diagnostic} 
invariant of the attention computation, not a 
\textbf{mechanistic} one. Formally:
\begin{enumerate}[nosep,label=(\roman*)]
  \item $\lambda_1$ is a deterministic function of the 
    attention matrix: $\lambda_1 = f(A^{(l,h)})$.
  \item Modifying $\lambda_1$ while preserving $A$ is 
    impossible; modifying $A$ to change $\lambda_1$ 
    also changes the entire spectral structure.
  \item The Fiedler partition targets the \emph{minimum cut}---
    the tokens already most weakly connected. Reducing their 
    residual connection changes nothing because those 
    attention weights were already near zero.
\end{enumerate}
\end{theorem}

\begin{remark}[Why this strengthens the retrieval claim]
\label{rem:strengthens}
The RED causal result is not a failure of the spectral 
retrieval theory---it is a \emph{confirmation} of its 
structure. The Fiedler partition reads out which tokens 
the head has already decided are relevant. You cannot 
change the decision by modifying the readout, any more 
than you can change the temperature by painting the 
thermometer. But the thermometer reading is still 
\emph{informative}: it tells you something true about 
the system that you cannot learn from the final hidden 
state alone. The claim is not ``$\lambda_1$ controls 
generation'' but ``$\phi_1$ reveals constraint structure 
that the residual stream compresses away.''
\end{remark}


%====================================================
\section{Verification Protocol: Spectral Information Recovery}
\label{sec:protocol}
%====================================================

\begin{definition}[Protocol 11.5: Spectral Information Recovery]
\label{def:protocol}
\textbf{Model:} GPT-2 family (124M--1.5B parameters).

\textbf{Task:} Factual completion with verifiable ground truth 
(e.g., ``The capital of France is'' $\to$ ``Paris'').

\textbf{Procedure:}
\begin{enumerate}[nosep]
  \item \textbf{Collect failure cases.} Run the model on $N$ 
    factual prompts. Identify prompts where 
    $\arg\max W_u h_L \neq \sigma^*$ (the correct answer).
  \item \textbf{Extract spectral structure.} For each failure 
    case, compute $\phi_1^{(l,h)}$ and $\lambda_1^{(l,h)}$ at 
    all $L \times H$ heads at the final token position.
  \item \textbf{Identify high-capacity heads.} Select heads 
    where $\lambda_1^{(l,h)} > \lambda_{\text{threshold}}$ 
    (determined from the empirically discriminating heads: 
    L3\,H8, L6\,H0, L3\,H7, L3\,H2).
  \item \textbf{Check partition alignment.} For each 
    high-capacity head, determine whether the correct-answer 
    token(s) appearing earlier in the prompt are in $V^+$ 
    (same side as the query position) while distractor tokens 
    are in $V^-$.
  \item \textbf{Spectral vote.} Aggregate across heads: if a 
    majority of high-capacity heads place the correct answer 
    in $V^+$, the spectral readout ``recovers'' the answer.
\end{enumerate}

\textbf{Pass criterion:} On the set of standard-readout 
failures, the spectral vote recovers the correct answer at 
rate $r > r_{\text{chance}}$ with $p < 0.01$. The information 
gap $\Delta S = S_{\text{spectral}} - S_{\text{hidden}}$ is 
measured as the log-odds ratio of spectral recovery rate vs.\ 
the model's own top-1 accuracy on those prompts (which is 0 by 
construction).

\textbf{Strong pass:} $r > 0.3$ (spectral readout recovers 
$> 30\%$ of failures). This would demonstrate that the 
constraint localization is computed correctly but lost in 
the residual-stream compression.
\end{definition}


%====================================================
\section{Synthesis: The Complete Chain}
\label{sec:synthesis}
%====================================================

We can now state the complete chain of identifications linking 
the abstract framework to the computable attention structure:

\begin{center}
\begin{tabular}{lll}
\toprule
\textbf{Abstract (Framework)} & \textbf{Spectral (Attention)} 
  & \textbf{Reference} \\
\midrule
Semantic manifold $\Mman$ & Token graph $G^{(l,h)}$ 
  at head $(l,h)$ & Def.~\ref{def:attn_graph} \\
Curvature $K$ & Spectral gap $\lambda_1^{(l,h)}$ 
  & Thm.~\ref{thm:gap_curvature} \\
Constraint localization $\psi$ & Fiedler partition 
  $\sign(\phi_1^{(l,h)})$ & Thm.~\ref{thm:fiedler_psi} \\
Davis Law $C = \tau/K$ & Retrieval capacity 
  $C^{(l,h)} = \tau / \lambda_1^{(l,h)}$ 
  & Eq.~\eqref{eq:C_head} \\
Trichotomy $\Gamma \gtrless 1$ & Partition sharpness 
  $\lambda_1 \gtrless \tau$ & Prop.~\ref{prop:trichotomy} \\
Abstention ($\lambda_1 < \lambda_{\min}$) & No clean 
  partition; head diffuses & \cite{davis2026geodesic} 
  Thm.~8.1 \\
Spectral capacity $\Csp = \lambda_1 D^2$ & Graph 
  conductance $\times$ diameter$^2$ 
  & \cite{davis2026spectral} Thm.~3.1 \\
Fisher metric $g^F$ & Attention-weighted token 
  similarity & \cite{davis2026info} Thm.~2.1 \\
\bottomrule
\end{tabular}
\end{center}

\medskip

The column on the left is the theory. The column in the 
middle is what we can compute from a single forward pass. 
Every abstract object has a concrete realization in the 
attention matrices.

The original framework described a \emph{Geometric Recurrent 
Neural Network} that navigates a manifold using $C = \tau/K$ 
to constrain its search. What the spectral retrieval 
identification shows is that standard transformers are 
\emph{already doing this}---each attention head is already 
computing a constraint localization via its Fiedler partition, 
already measuring curvature via its spectral gap, already 
implementing the Davis Law. They just read out the result 
through a lossy linear projection (the residual stream and 
unembedding matrix) instead of through the spectral structure 
directly.

The question is no longer ``Can we build a geometric NN?'' 
It is: ``Can we read the geometry that transformers already 
compute?''


%====================================================
% BIBLIOGRAPHY
%====================================================
\begin{thebibliography}{99}

\bibitem{davis2026geodesic}
Davis, B.~R. (2026). Geodesic Computation: Mathematical 
Foundations for Geometry-Native Intelligence. Zenodo. 
doi:10.5281/zenodo.18673121.

\bibitem{davis2026field}
Davis, B.~R. (2026). The Field Equations of Semantic 
Coherence. Unpublished manuscript.

\bibitem{davis2026spectral}
Davis, B.~R. (2026). Branch~IX: Spectral Geometry of 
Completion. Unpublished manuscript.

\bibitem{davis2026info}
Davis, B.~R. (2026). Branch~X: Information Geometry of 
Semantic Fields. Unpublished manuscript.

\bibitem{chung1997}
Chung, F.~R.~K. (1997). \emph{Spectral Graph Theory}. 
AMS.

\bibitem{fiedler1973}
Fiedler, M. (1973). Algebraic connectivity of graphs. 
\emph{Czechoslovak Mathematical Journal}, 23(2), 298--305.

\bibitem{cheeger1970}
Cheeger, J. (1970). A lower bound for the smallest 
eigenvalue of the Laplacian. In \emph{Problems in 
Analysis} (pp.~195--199). Princeton Univ.\ Press.

\end{thebibliography}

\end{document}
